\chapter{Lecture Notes}

\section{Lecture 1}

\begin{convention}
Rings are commutative with $1$.
\end{convention}

\begin{definition}{}{def:ring-morphism}
Let $R, S$ be rings. A \emph{morphism} $\varphi$ from $R$ to $S$ is a map $\varphi: R \to S$ such that
\begin{enumerate}[label=(\arabic*)]
    \item $\varphi(r + s) = \varphi(r) + \varphi(s)$,
    \item $\varphi(rs) = \varphi(r)\varphi(s)$,
    \item $\varphi(1) = 1$.
\end{enumerate}
\end{definition}

\begin{convention}
A subring $R$ must contain $1$.
\end{convention}

\begin{question}
What is $R[X]$?
\end{question}

$X$ is not an element of $R$. Consider
\[
S = \left\{ f : \bbn \to R \ \middle|\ f(r) = 0 \text{ for all but finitely many } r \right\},
\]
with the usual addition and multiplication in sequences (Cauchy product). Then $S$ is a ring with identity
\[
1_S = (1, 0, 0, \ldots, 0).
\]

Define $\varphi: R \to S$ by $\varphi(a) = (a, 0, \ldots, 0)$. Then $\varphi$ is a ring homomorphism, for sure.

Define $X = (0, 1, 0, \ldots)$ and observe $X^2 = (0, 0, 1, 0, \ldots)$. A polynomial $P$ over $R$ can then be identified as an element of $S$, say
\[
P = (P_0, P_1, \ldots, P_n, 0, \ldots) = \varphi(P_0) \cdot 1_S + \varphi(P_1) \cdot X + \cdots + \varphi(P_n) \cdot X^n.
\]

\begin{definition}{}{def:prime-element}
Let $R$ be a ring. A non-zero, non-unit element $p \in R$ is \emph{prime} if
\[
p \mid ab \implies p \mid a \text{ or } p \mid b
\]
(include the elementwise definition: for $a, b \in R$, $a \mid b$ if $\exists\, r \in R$ such that $b = r \cdot a$).
\end{definition}

\begin{definition}{}{def:prime-ideal}
Let $R$ be a ring. A proper ideal $P$ of $R$ is called a \emph{prime ideal} if
\[
ab \in P \implies a \in P \text{ or } b \in P.
\]
\end{definition}

\begin{definition}{}{def:ideal}
Let $R$ be a ring. $I \subseteq R$ is called an \emph{ideal} if there exists a ring $S$ and a morphism $\varphi: R \to S$ such that $I = \ker(\varphi)$.
\end{definition}

\begin{remark}{}{rem:ideal-elementwise}
This immediately gives the elementwise definition of ideals.
\end{remark}

\[
\caln(R) = \{\text{all nilpotent elements of } R\} \quad \leadsto \quad \text{the \emph{nil-radical} of } R.
\]

\begin{theorem}{}{thm:nilradical}
\[
\caln(R) = \bigcap_{P \in \Spec R} P, \qquad \text{where } \Spec R = \text{all prime ideals of } R.
\]
\end{theorem}

\begin{proof}
Let $x \in \caln(R)$. Then for any $P \in \Spec R$, there exists $n_P \in \bbn$ such that
\[
x^{n_P} = 0 \implies 0 \in P \implies x^{n_P} \in P \implies x \in P.
\]

Now let $x \notin \caln(R)$, and define
\[
S = \{ I \subseteq R \text{ ideal} \mid x^n \notin I \text{ for all } n \}.
\]
$S \neq \emptyset$, as $\langle 0 \rangle \in S$, and $S$ is a poset (ordered by inclusion). By Zorn's Lemma, let $M$ be a maximal element of $S$.

\textbf{Claim.} $M$ is a prime ideal.

\emph{Proof of claim.} Let $a \notin M$, $b \notin M$. Now,
\[
M + \langle a \rangle, \ M + \langle b \rangle \notin S \implies \exists\, m, n \in \bbn \text{ such that}
\]
\[
x^m \in M + \langle a \rangle \text{ and } x^n \in M + \langle b \rangle \implies x^{m+n} \in M + \langle ab \rangle.
\]
If $ab \in M$, then $x^{m+n} \in M$, a contradiction. Hence $ab \notin M$, so $M$ is prime. $\hfill\square$

Since $x \notin M$ and $M \in \Spec R$, we conclude $x \notin \bigcap_{P \in \Spec R} P$, completing the proof.
\end{proof}

\begin{question}
What are the units of $R[X]$?
\end{question}

If $R$ is an integral domain, then they are precisely the units of $R$.

\begin{proposition}{}{prop:units-poly}
$f(X) = a_n X^n + \cdots + a_1 X + a_0$ is a unit if $a_0$ is a unit and the $a_i$'s are nilpotent.
\end{proposition}

\begin{proof}
Let $P \in \Spec R$. We have a natural morphism $\varphi: R[X] \to (R/P)[X]$ given by
\[
\varphi(b_m X^m + \cdots + b_1 X + b_0) = \overline{b_m} X^m + \cdots + \overline{b_1} X + \overline{b_0}, \qquad (\overline{b_i} = b_i \bmod P).
\]
Then $\varphi(f)$ is a unit in $(R/P)[X]$, and since $R/P$ is an integral domain, this forces
\[
\overline{a_i} \equiv \overline{0} \pmod{P} \text{ for all } i \in [n], \qquad \text{and} \qquad \overline{a_0} \text{ is a unit of } R/P.
\]
That is, $a_0 \notin P$ and $a_1, \ldots, a_n \in P$. As $P \in \Spec R$ was arbitrary,
\[
a_0 \notin P \text{ and } a_i \in P \ \ \forall\, i \in [n], \qquad \text{for all } P \in \Spec R.
\]

The other direction is trivial.
\end{proof}

\section{Lecture 2}

\begin{definition}{}{def:module}
Let $R$ be a ring and let $(M, +)$ be an abelian group. A \emph{scalar multiplication}
\[
R \times M \longrightarrow M, \qquad (r, m) \longmapsto rm,
\]
makes $M$ into an \emph{$R$-module} if the following hold for all $r, s \in R$ and $m, m_1, m_2 \in M$:
\begin{enumerate}[label=(\roman*)]
    \item $(r + s)(m) = rm + sm$,
    \item $r(m_1 + m_2) = rm_1 + rm_2$,
    \item $r(sm) = rsm = s(rm)$,
    \item $1 \cdot m = m$.
\end{enumerate}
\end{definition}

\begin{definition}{}{def:module-morphism}
If $M, N$ are $R$-modules, a map $\varphi: M \to N$ is called an \emph{$R$-module morphism} if
\begin{enumerate}[label=(\roman*)]
    \item $\varphi(m_1 + m_2) = \varphi(m_1) + \varphi(m_2)$,
    \item $\varphi(rm) = r\varphi(m)$.
\end{enumerate}
\end{definition}

\begin{definition}{}{def:free-module}
Let $R$ be a ring and $S$ a nonempty set. A \emph{free $R$-module on $S$} is a pair $(F, f)$, where $F$ is an $R$-module and $f: S \to F$ is a set-theoretic map, such that: for any $R$-module $M$ and any set-theoretic map $g: S \to M$, there exists a unique $R$-module morphism $h: F \to M$ such that
\[
\begin{tikzcd}
S \arrow[r, "g"] \arrow[d, "f"'] & M \\
F \arrow[ur, dashed, "h"'] &
\end{tikzcd}
\]
commutes, i.e.\ $h \circ f = g$.
\end{definition}

\begin{proposition}{}{prop:free-module-properties}
Let $(F, f)$ be a free $R$-module on $S$. Then:
\begin{enumerate}[label=(\arabic*)]
    \item $f$ is injective;
    \item $\Im(f)$ generates $F$ as an $R$-module.
\end{enumerate}
\end{proposition}

\begin{proof}[Proof of (1)]
Let $x, y \in S$ with $x \neq y$. \textbf{Goal:} $f(x) \neq f(y)$.

Take any $R$-module $M$ with at least two elements, and define $g: S \to M$ such that $g(x) \neq g(y)$. By definition, there exists a (unique) morphism $h: F \to M$ such that $h \circ f = g$. Then
\[
(h \circ f)(x) = g(x) \neq g(y) = (h \circ f)(y) \implies f(x) \neq f(y). \qedhere
\]
\end{proof}

\begin{definition}{}{def:submodule-generated}
Let $M$ be an $R$-module and $\emptyset \neq A \subseteq M$. The \emph{submodule of $M$ generated by $A$} is, equivalently:
\begin{enumerate}[label=(\arabic*)]
    \item the smallest submodule of $M$ containing $A$;
    \item the intersection of all submodules of $M$ containing $A$;
    \item the set of all finite $R$-linear combinations of elements of $A$.
\end{enumerate}
\end{definition}

\begin{proof}[Proof of (2)]
Let $A$ be the submodule of $F$ generated by $\Im(f)$. Since $f(s) \in \Im(f) \subseteq A$ for every $s \in S$, $f$ factors as $f = i_A \circ f^*$ for a set-theoretic map $f^*: S \to A$ and the inclusion $i_A: A \hookrightarrow F$:
\[
\begin{tikzcd}
S \arrow[r, "f^*"] \arrow[dr, "f"'] & A \arrow[d, hook, "i_A"] \\
& F
\end{tikzcd}
\]

As $(F, f)$ is free on $S$ and $A$ is an $R$-module, there is a unique morphism $h: F \to A$ such that $h \circ f = f^*$:
\[
\begin{tikzcd}
S \arrow[r, "f^*"] \arrow[d, "f"'] & A \\
F \arrow[ur, dashed, "h"'] &
\end{tikzcd}
\]

It is enough to show $i_A$ is onto. We compute
\[
(i_A \circ h) \circ f = i_A \circ (h \circ f) = i_A \circ f^* = f.
\]
So $i_A \circ h$ and $\mathrm{id}_F$ both satisfy $(-) \circ f = f$; by uniqueness in the definition of freeness (applied with $M = F$, $g = f$), $i_A \circ h = \mathrm{id}_F$. Hence $i_A$ is onto, i.e.\ $A = F$, so $\Im(f)$ generates $F$.
\end{proof}

\begin{theorem}{Uniqueness of free modules}{thm:free-uniqueness}
Let $(F, f)$ be a free $R$-module on a set $S$. Then $(F', f')$ is also a free $R$-module on $S$ if and only if there exists a unique $R$-module isomorphism $j: F \to F'$ such that $j \circ f = f'$.
\end{theorem}

\begin{proof}
Suppose $(F', f')$ is also free on $S$. Since $(F, f)$ is free (applied to $M = F'$, $g = f'$), there is a morphism $h: F \to F'$ with $h \circ f = f'$:
\[
\begin{tikzcd}
S \arrow[r, "f'"] \arrow[d, "f"'] & F' \\
F \arrow[ur, "h"'] &
\end{tikzcd}
\]
Likewise, since $(F', f')$ is free (applied to $M = F$, $g = f$), there is a morphism $h': F' \to F$ with $h' \circ f' = f$:
\[
\begin{tikzcd}
S \arrow[r, "f"] \arrow[d, "f'"'] & F \\
F' \arrow[ur, "h'"'] &
\end{tikzcd}
\]
Then $(h' \circ h) \circ f = h' \circ (h \circ f) = h' \circ f' = f$. So $h' \circ h$ and $\mathrm{id}_F$ both satisfy $(-) \circ f = f$; by uniqueness (with $M = F$, $g = f$), $h' \circ h = \mathrm{id}_F$. Similarly $h \circ h' = \mathrm{id}_{F'}$. So $j := h$ is an isomorphism with $j \circ f = f'$.

\textbf{Converse.} Suppose there is an isomorphism $j: F \to F'$ with $j \circ f = f'$. To show $(F', f')$ is free on $S$, take any $R$-module $M$ and set-theoretic map $g: S \to M$; we must produce a unique $h': F' \to M$ with $h' \circ f' = g$:
\[
\begin{tikzcd}
S \arrow[r, "g"] \arrow[d, "f'"'] & M \\
F' \arrow[ur, dashed, "h'"'] &
\end{tikzcd}
\]
Since $(F,f)$ is free, there is a unique $h: F \to M$ with $h \circ f = g$. From $j \circ f = f'$ we get $j^{-1} \circ f' = f$, so
\[
(h \circ j^{-1}) \circ f' = h \circ (j^{-1} \circ f') = h \circ f = g,
\]
i.e.\ $h' := h \circ j^{-1}$ works:
\[
\begin{tikzcd}
S \arrow[r, "g"] \arrow[d, "f'"'] \arrow[dr, "f"'] & M \\
F' \arrow[r, "j^{-1}"'] & F \arrow[u, "h"']
\end{tikzcd}
\]

\textbf{Uniqueness.} Suppose $t: F' \to M$ also satisfies $t \circ f' = g$. Then $t \circ j: F \to M$ satisfies $(t \circ j) \circ f = t \circ (j \circ f) = t \circ f' = g$, so by uniqueness of $h$ (from freeness of $(F,f)$), $t \circ j = h$, hence $t = h \circ j^{-1} = h'$.
\end{proof}

\begin{theorem}{Existence of free $R$-modules}{thm:free-existence}
For any ring $R$ and nonempty set $S$, a free $R$-module on $S$ exists.
\end{theorem}

\begin{proof}
Define
\[
F = \{ \theta: S \to R \mid \theta(s) = 0 \text{ for almost all } s \in S \}.
\]
For $\theta_1, \theta_2 \in F$ and $\lambda \in R$, define
\[
(\theta_1 + \theta_2)(s) = \theta_1(s) + \theta_2(s), \qquad (\lambda \theta)(s) = \lambda\, \theta(s).
\]
One checks that $F$ is an $R$-module. Define $f: S \to F$ by $f(s) = \delta_s$, where $\delta_s(t) = 1$ if $t = s$ and $0$ otherwise.

\textbf{Existence of $h$.} Let $M$ be an $R$-module and $g: S \to M$ a set-theoretic map. Define $h: F \to M$ by
\[
h(\theta) = \sum_{s \in S} \theta(s)\, g(s),
\]
a finite sum for every $\theta \in F$. For any $t \in S$,
\[
(h \circ f)(t) = h(f(t)) = h(\delta_t) = \delta_t(t)\, g(t) = g(t),
\]
so $h \circ f = g$:
\[
\begin{tikzcd}
S \arrow[r, "g"] \arrow[d, "f"'] & M \\
F \arrow[ur, dashed, "h"'] &
\end{tikzcd}
\]

\textbf{Uniqueness of $h$.} For $\theta \in F$ and $t \in S$,
\[
\theta(t) = \theta(t) \cdot 1_R = \sum_{s \in S} \theta(s)\, \delta_s(t) = \left( \sum_{s \in S} \theta(s)\, \delta_s \right)(t) \qquad \forall\, t \in S,
\]
so $\theta = \sum_{s \in S} \theta(s)\, \delta_s$. Now let $h': F \to M$ be any $R$-module morphism with $h' \circ f = g$. Then for $\theta \in F$,
\[
h'(\theta) = h'\left( \sum_{s \in S} \theta(s)\, \delta_s \right) = \sum_{s \in S} \theta(s)\, h'(\delta_s) = \sum_{s \in S} \theta(s)\, h'(f(s)) = \sum_{s \in S} \theta(s)\, g(s) = h(\theta).
\]
Hence $h = h'$.
\end{proof}

\section{Lecture 3}

\begin{theorem}{}{thm:prime-implies-principal}
Let $R$ be a ring such that every prime ideal is principal. Then every ideal is principal.
\end{theorem}

\begin{proof}
Let $S = \{\text{all non-principal ideals of } R\}$. For the sake of contradiction (FTSOC), suppose $S \neq \emptyset$. By Zorn's Lemma, let $I$ be a maximal element of $S$.

Since $I$ is not principal, $I$ is not prime by hypothesis, so there exist $a, b \in R$ with $a \notin I$, $b \notin I$, but $ab \in I$.

Let $J = \langle a \rangle + I \supsetneq I$; since $I$ is maximal among non-principal ideals, $J \notin S$, so $J$ is principal, say $J = \langle c \rangle$.

Also let $K = \{ r \in R \mid ra \in I \} \supsetneq I$ (since $b \in K \setminus I$); again $K \notin S$, so $K = \langle d \rangle$ for some $d$.

\textbf{$I \subseteq \langle cd \rangle$.} Let $x \in I$. Since $I \subseteq J = \langle c \rangle$, there is $\lambda \in R$ with $\lambda c = x$. Now $\lambda a \in I$, so $\lambda \in K = \langle d \rangle$, and hence $x \in \langle cd \rangle$. Thus $I \subseteq \langle cd \rangle$.

\textbf{$\langle cd \rangle \subseteq I$.} There exist $u \in R$, $v \in I$ such that $c = ua + v$ (since $c$ generates $J = \langle a \rangle + I$). Also $da \in I$ (since $d$ generates $K$). Then
\[
cd = (ua+v)d = u(ad) + vd \in I,
\]
since $ad = da \in I$ and $v \in I$. Hence $\langle cd \rangle \subseteq I$.

Together, $I = \langle cd \rangle$, contradicting $I \in S$. So $S = \emptyset$, i.e.\ every ideal of $R$ is principal.
\end{proof}

\begin{theorem}{}{thm:cohen}
Let $R$ be a ring such that every prime ideal is finitely generated. Then every ideal of $R$ is finitely generated.
\end{theorem}

\begin{proof}
Let $A = \{ I \text{ ideal of } R \mid I \text{ is not finitely generated} \}$. By Zorn's Lemma, let $M$ be a maximal element of $A$.

\textbf{Claim: $M$ is a maximal ideal.} Suppose not, so there is an ideal $N$ with $M \subsetneq N \subsetneq R$. Since $N$ is maximal (hence prime), it is finitely generated by hypothesis: there exist $a_1, \ldots, a_n \in R$ such that $N = \langle a_1, \ldots, a_n \rangle$.

Since $M \subsetneq N$, some $a_i \notin M$. Then $M + \langle a_i \rangle \supsetneq M$, so by maximality of $M$ in $A$, $M + \langle a_i \rangle \notin A$; hence there exist $b_1, \ldots, b_k \in R$ such that
\[
M + \langle a_i \rangle = \langle b_1, \ldots, b_k \rangle.
\]
\end{proof}

\begin{lemma}{}{lem:union-two-ideals}
Let $I, J_1, J_2$ be ideals of $R$ such that $I \subseteq J_1 \cup J_2$. Then either $I \subseteq J_1$ or $I \subseteq J_2$.
\end{lemma}

\begin{proof}
FTSOC, suppose $I \setminus J_1 \neq \emptyset$ and $I \setminus J_2 \neq \emptyset$. Let $x \in I \setminus J_1$ and $y \in I \setminus J_2$; since $I \subseteq J_1 \cup J_2$, this forces $x \in J_2$ and $y \in J_1$. Also $x+y \in I \subseteq J_1 \cup J_2$, so $x+y \in J_1$ or $x+y \in J_2$; WLOG $x+y \in J_1$. Since $y \in J_1$, $x = (x+y) - y \in J_1$, contradicting $x \notin J_1$.
\end{proof}

\begin{theorem}{Prime Avoidance}{thm:prime-avoidance}
Let $I, J_1, \ldots, J_n$ be ideals of $R$ such that
\begin{enumerate}[label=(\arabic*)]
    \item $I \subseteq J_1 \cup \cdots \cup J_n$,
    \item $J_3, \ldots, J_n \in \Spec R$.
\end{enumerate}
Then $I \subseteq J_i$ for some $i$.
\end{theorem}

\begin{proof}
We induct on $n$. The case $n=2$ needs no primality hypothesis and is the lemma above; assume $n \geq 3$ and the result for any collection of $n-1$ such ideals.

We prove the contrapositive: if $I \not\subseteq J_i$ for every $i$, then $I \not\subseteq \bigcup_{i=1}^n J_i$.

Fix $i$. Apply the induction hypothesis to $I$ together with the $n-1$ ideals $\{J_j\}_{j \neq i}$ (which still has all but at most two members prime): if $I \subseteq \bigcup_{j \neq i} J_j$, the induction hypothesis would give $I \subseteq J_j$ for some $j \neq i$, contradicting our assumption. So $I \not\subseteq \bigcup_{j \neq i} J_j$, and we may choose
\[
z_i \in I \setminus \bigcup_{j \neq i} J_j.
\]
Since $I \subseteq \bigcup_{i=1}^n J_i$ by hypothesis (1), and $z_i \notin J_j$ for every $j \neq i$, we must have $z_i \in J_i$, for every $i$.

Set
\[
z := z_1 z_2 \cdots z_{n-1} + z_n \in I
\]
(a sum of an element of $I$ and a product of elements of $I$, hence itself in $I$). We show $z \notin J_i$ for every $i$, which gives $I \not\subseteq \bigcup_{i=1}^n J_i$, as desired.

\begin{itemize}
    \item \emph{If $i \in \{1, \ldots, n-1\}$:} since $z_i \in J_i$ is one of the factors of $z_1 \cdots z_{n-1}$, this product lies in $J_i$. If $z \in J_i$, then $z_n = z - z_1 \cdots z_{n-1} \in J_i$, contradicting $z_n \notin J_i$ (as $i \neq n$).
    \item \emph{If $i = n$:} if $z \in J_n$, then since $z_n \in J_n$ as well, $z_1 \cdots z_{n-1} = z - z_n \in J_n$. As $J_n$ is prime, some factor $z_k \in J_n$ for a $k \in \{1, \ldots, n-1\}$, contradicting $z_k \notin J_n$ (as $k \neq n$).
\end{itemize}
This completes the induction.
\end{proof}

\begin{proposition}{}{prop:integral-closure-Zi}
An element of $\bbq(i)$ lies in $\bbz[i]$ if and only if it is the root of a monic polynomial over $\bbz$.
\end{proposition}

\begin{proof}
Let $a + bi \in \bbq(i)$. Then $a+bi$ is a root of
\[
P(X) = X^2 - 2aX + (a^2+b^2) \in \bbq[X].
\]

\textbf{Claim: $2a \in \bbz$ and $a^2+b^2 \in \bbz$ together imply $a, b \in \bbz$.}

Since $2a \in \bbz$, $(2a)^2 \in \bbz$. Since $a^2+b^2 \in \bbz$, $4(a^2+b^2) = (2a)^2 + (2b)^2 \in \bbz$, so $(2b)^2 = 4(a^2+b^2) - (2a)^2 \in \bbz$. As $2b \in \bbq$ has integer square, $2b \in \bbz$.

Now $(2a)^2 + (2b)^2 = 4(a^2+b^2) \equiv 0 \pmod 4$. Since the square of any integer is $\equiv 0 \pmod 4$ if it is even, and $\equiv 1 \pmod 4$ if it is odd, the only way for the sum to be $\equiv 0 \pmod 4$ is
\[
(2a)^2 \equiv (2b)^2 \equiv 0 \pmod 4,
\]
forcing $2a, 2b$ to both be even, i.e.\ $a, b \in \bbz$.

Since $a+bi \in \bbz[i]$ exactly when $a, b \in \bbz$, and $a+bi$ is a root of a monic polynomial over $\bbz$ exactly when $2a, a^2+b^2 \in \bbz$ (the coefficients of $P(X)$), the claim gives the equivalence.
\end{proof}

\begin{definition}{}{def:integral-element}
Let $R \hookrightarrow S$ be rings. An element $\alpha \in S$ is called \emph{integral} over $R$ if $\alpha$ is a root of a monic polynomial
\[
X^n + a_{n-1}X^{n-1} + \cdots + a_0
\]
over $R$ (i.e.\ with $a_0, \ldots, a_{n-1} \in R$).
\end{definition}